\chapter{2025年上海卷（春）}

\section{填空题}

\begin{problem}
已知集合 \(A=\{x\mid x>0\}\)，\(B=\{-1,0,1,2\}\)，则 \(A\cap B=\fillinblank{}\).
\end{problem}

\begin{answer}
\(\{1,2\}\)
\end{answer}

\begin{solution}
集合 \(A\) 中的元素需满足 \(x>0\). 在 \(B=\{-1,0,1,2\}\) 中满足该条件的元素是 \(1,2\)，故 \(A\cap B=\{1,2\}\).
\end{solution}

\begin{problem}
不等式 \(\dfrac{x}{x-1}<0\) 的解集为\fillinblank{}.
\end{problem}

\begin{answer}
\((0,1)\)
\end{answer}

\begin{solution}
分式小于零等价于分子、分母异号. 临界点为 \(0,1\)，在数轴上讨论可得当 \(0<x<1\) 时 \(\dfrac{x}{x-1}<0\)，故解集为 \((0,1)\).
\end{solution}

\begin{problem}
已知复数 \(z=\dfrac{2+\i}{\i}\)，其中 \(\i\) 为虚数单位，则 \(|z|=\fillinblank{}\).
\end{problem}

\begin{answer}
\(\sqrt5\)
\end{answer}

\begin{solution}
因为 \(z=\dfrac{2+\i}{\i}=(2+\i)(-\i)=1-2\i\)，所以 \(|z|=\sqrt{1^2+(-2)^2}=\sqrt5\).
\end{solution}

\begin{problem}
已知 \(\vec a=(2,1)\)，\(\vec b=(1,x)\)，若 \(\vec a\parallel\vec b\)，则 \(x=\fillinblank{}\).
\end{problem}

\begin{answer}
\(\dfrac12\)
\end{answer}

\begin{solution}
由 \(\vec a\parallel\vec b\) 得 \(2x-1\cdot1=0\)，故 \(x=\dfrac12\).
\end{solution}

\begin{problem}
已知 \(\tan\alpha=1\)，则 \(\cos\left(\alpha+\dfrac{\pi}{4}\right)=\fillinblank{}\).
\end{problem}

\begin{answer}
\(0\)
\end{answer}

\begin{solution}
由 \(\tan\alpha=1\) 得 \(\sin\alpha=\cos\alpha\). 因此 \(\cos\left(\alpha+\dfrac{\pi}{4}\right)=\dfrac{\sqrt2}{2}(\cos\alpha-\sin\alpha)=0\).
\end{solution}

\begin{problem}
已知 \(\left(x+\dfrac{m}{x}\right)^6\) 的展开式中常数项为 \(20\)，则实数 \(m\) 的值为\fillinblank{}.
\end{problem}

\begin{answer}
\(1\)
\end{answer}

\begin{solution}
通项为 \(\binom6k x^{6-k}\left(\dfrac{m}{x}\right)^k=\binom6k m^k x^{6-2k}\). 常数项满足 \(6-2k=0\)，即 \(k=3\). 故常数项为 \(\binom63m^3=20m^3=20\)，得 \(m=1\).
\end{solution}

\begin{problem}
已知 \(\{a_n\}\) 是首项为 \(1\)、公差为 \(1\) 的等差数列，\(\{b_n\}\) 是首项为 \(1\)、公比为 \(q(q>0)\) 的等比数列. 若数列 \(\{a_nb_n\}\) 的前三项和为 \(2\)，则 \(q=\fillinblank{}\).
\end{problem}

\begin{answer}
\(\dfrac13\)
\end{answer}

\begin{solution}
由题意 \(a_1=1,a_2=2,a_3=3\)，且 \(b_1=1,b_2=q,b_3=q^2\). 因此前三项和为 \(1+2q+3q^2=2\)，即 \(3q^2+2q-1=0\). 因 \(q>0\)，得 \(q=\dfrac13\).
\end{solution}

\begin{problem}
关于 \(x\) 的方程 \(|x-1|+|\pi-x|=\pi-1\) 的解集为\fillinblank{}.
\end{problem}

\begin{answer}
\([1,\pi]\)
\end{answer}

\begin{solution}
当 \(1\le x\le\pi\) 时，\(|x-1|+|\pi-x|=x-1+\pi-x=\pi-1\). 当 \(x<1\) 或 \(x>\pi\) 时，两点间距离和大于 \(\pi-1\). 故解集为 \([1,\pi]\).
\end{solution}

\begin{problem}
已知 \(P\) 是一个圆锥的顶点，\(PA\) 是母线，\(PA=2\)，该圆锥的底面半径是 \(1\). \(B,C\) 分别在圆锥的底面上，则异面直线 \(PA\) 与 \(BC\) 所成角的最小值为\fillinblank{}.
\end{problem}

\begin{answer}
\(\dfrac{\pi}{3}\)
\end{answer}

\begin{solution}
圆锥轴截面中，母线长为 \(2\)，底面半径为 \(1\)，故母线与底面所成角满足 \(\cos\theta=\dfrac{1}{2}\)，即 \(\theta=\dfrac{\pi}{3}\). 底面弦 \(BC\) 的方向可在底面内变化，异面直线所成角的最小值为母线与底面内直线所成角的最小值，故为 \(\dfrac{\pi}{3}\).
\end{solution}

\begin{problem}
已知双曲线 \(\dfrac{x^2}{a^2}-\dfrac{y^2}{6-a^2}=1(a>0)\) 的左、右焦点分别为 \(F_1,F_2\). 通过 \(F_2\) 且倾斜角为 \(\dfrac{\pi}{3}\) 的直线与双曲线交于第一象限的点 \(A\)，延长 \(AF_2\) 至 \(B\) 使得 \(AB=AF_1\). 若 \(\triangle BF_1F_2\) 的面积为 \(3\sqrt6\)，则 \(a=\fillinblank{}\).
\end{problem}

\begin{answer}
\(\sqrt3\)
\end{answer}

\begin{solution}
该双曲线中 \(c^2=a^2+(6-a^2)=6\)，故 \(|F_1F_2|=2\sqrt6\). 直线 \(AF_2\) 的倾斜角为 \(\dfrac{\pi}{3}\)，点 \(B\) 到 \(x\) 轴的距离为 \(|BF_2|\sin\dfrac{\pi}{3}\). 由面积条件可得 \( \dfrac12\cdot2\sqrt6\cdot |BF_2|\sin\dfrac{\pi}{3}=3\sqrt6\)，从而 \(|BF_2|=2\sqrt3\). 又由双曲线定义与 \(AB=AF_1\)，可推出 \(|BF_2|=2a\)，故 \(a=\sqrt3\).
\end{solution}

\begin{problem}
如图，正方形 \(ABCD\) 是一块边长为 \(4\) 的工程用料，阴影部分是被腐蚀的区域，其余部分完好，曲线 \(MN\) 为以 \(AC\) 为对称轴的抛物线的一部分，且 \(DM=DN=3\). 从完好的部分中截取一块矩形原料，当其面积有最大值时，矩形一边的长为\fillinblank{}.

\bitmapfigure[max width=0.35\linewidth,max totalheight=4.2cm]{img/2025/shanghai_spring/q11_fig1.png}
\end{problem}

\begin{answer}
\(\dfrac{4+\sqrt7}{3}\)
\end{answer}

\begin{solution}
以 \(A\) 为原点建立坐标系，抛物线关于对角线对称. 设矩形相关边长为 \(x\)，由抛物线方程可把矩形面积化为三次函数，对其求导得驻点 \(x=\dfrac{4\pm\sqrt7}{3}\). 结合取值范围与面积最大条件，取 \(x=\dfrac{4+\sqrt7}{3}\).
\end{solution}

\begin{problem}
在平面中，\(\vec e_1\) 和 \(\vec e_2\) 是互相垂直的单位向量，向量 \(\vec a\) 满足 \(|\vec a-4\vec e_1|=2\)，向量 \(\vec b\) 满足 \(|\vec b-6\vec e_2|=1\)，求 \(\vec b\) 在 \(\vec a\) 方向上的数量投影的最大值\fillinblank{}.
\end{problem}

\begin{answer}
\(4\)
\end{answer}

\begin{solution}
把 \(\vec e_1,\vec e_2\) 作为坐标基底，则 \(\vec a\) 的端点在圆心 \((4,0)\)、半径 \(2\) 的圆上，\(\vec b\) 的端点在圆心 \((0,6)\)、半径 \(1\) 的圆上. \(\vec b\) 在 \(\vec a\) 方向上的数量投影为 \(\vec b\cdot \dfrac{\vec a}{|\vec a|}\). 沿可变方向取最大时，结合两圆的支撑函数计算可得最大值为 \(4\).
\end{solution}
\section{单选题}

\begin{problem}
如图，\(ABCD-A_1B_1C_1D_1\) 是正四棱台，则下列各组直线中属于异面直线的是（\quad）.

\bitmapfigure[max width=0.45\linewidth,max totalheight=4.0cm]{img/2025/shanghai_spring/q13_fig1.png}
\choices
    {\(AB\) 和 \(C_1D_1\)}
    {\(AA_1\) 和 \(CC_1\)}
    {\(BD_1\) 和 \(B_1D\)}
    {\(A_1D_1\) 和 \(AB\)}
\end{problem}

\begin{answer}
D
\end{answer}

\begin{solution}
\(AB\) 与 \(C_1D_1\) 平行，\(AA_1\) 与 \(CC_1\) 可共面，\(BD_1\) 与 \(B_1D\) 在同一对角截面内. \(A_1D_1\) 与 \(AB\) 既不平行也不相交，且不共面，故为异面直线.
\end{solution}

\begin{problem}
幂函数 \(y=x^\alpha\) 在 \((0,+\infty)\) 上是严格减函数，且经过 \((-1,-1)\)，则 \(\alpha\) 的值可能是（\quad）.
\choices
    {\(-\dfrac23\)}
    {\(-\dfrac13\)}
    {\(\dfrac13\)}
    {\(3\)}
\end{problem}

\begin{answer}
B
\end{answer}

\begin{solution}
幂函数在 \((0,+\infty)\) 上严格递减要求 \(\alpha<0\). 又图像经过 \((-1,-1)\)，指数需使负数处有定义且函数值为负. 选项中 \(\alpha=-\dfrac13\) 满足，故选 \(B\).
\end{solution}

\begin{problem}
有一四边形 \(ABCD\)，对其四边 \(AB,BC,CD,DA\) 按顺序分别抛掷一枚质量均匀的硬币：正面朝上则擦去该边，反面朝上则不擦去. 最后，以 \(A\) 为起点沿尚未擦去的边出发，可以到达 \(C\) 点的概率为（\quad）.
\choices
    {\(\dfrac12\)}
    {\(\dfrac7{16}\)}
    {\(\dfrac14\)}
    {\(\dfrac3{16}\)}
\end{problem}

\begin{answer}
B
\end{answer}

\begin{solution}
四条边各自保留或擦去，共 \(2^4=16\) 种等可能结果. 从 \(A\) 到 \(C\) 可经路径 \(A-B-C\) 或 \(A-D-C\). 不能到达 \(C\) 当且仅当两条路径都至少缺一条边. 枚举可得可达情形共 \(7\) 种，概率为 \(\dfrac7{16}\).
\end{solution}

\begin{problem}
已知 \(a\in\mathbb R\)，不等式 \(\left(\tan\dfrac{\pi x}{6}-a\right)\left(\tan\dfrac{\pi x}{6}-a-1\right)<0\) 在 \((0,2025)\) 中的整数解有 \(m\) 个. 关于 \(m\)，以下不可能的是（\quad）.
\choices
    {\(0\)}
    {\(338\)}
    {\(674\)}
    {\(1012\)}
\end{problem}

\begin{answer}
D
\end{answer}

\begin{solution}
函数 \(\tan\dfrac{\pi x}{6}\) 的周期为 \(6\). 在整数点上，其取值按周期重复，并且每个周期中至多有若干固定位置落入开区间 \((a,a+1)\). 区间 \((0,2025)\) 含 \(337\) 个完整周期再加余项，因此可出现的整数解个数受周期结构限制，不可能达到 \(1012\). 故选 \(D\).
\end{solution}
\section{解答题}

\begin{problem}
在三棱锥 \(P-ABC\) 中，平面 \(PAB\perp\) 平面 \(ABC\)，\(PA=AB=CB=2\)，\(AC=BC=\sqrt2\).
\begin{enumerate}
\item 若 \(O\) 是棱 \(AC\) 的中点，证明：\(BO\perp\) 平面 \(PAC\)，并求三棱锥 \(B-OCP\) 的体积；
\item 求二面角 \(B-PC-A\) 的大小.
\end{enumerate}
\end{problem}

\begin{solution}
（1） 因为 \(AC=BC=\sqrt2\)，\(AB=2\)，所以 \(\triangle ABC\) 是以 \(C\) 为直角的等腰直角三角形，\(O\) 为 \(AC\) 中点时可由中点与垂直关系推出 \(BO\perp AC\). 又平面 \(PAB\perp\) 平面 \(ABC\)，结合 \(PA=AB\) 的对称关系可得 \(BO\perp PO\)，故 \(BO\perp\) 平面 \(PAC\). 以 \(\triangle OCP\) 为底、\(BO\) 为高，计算得三棱锥体积为 \(\dfrac{\sqrt3}{6}\).

（2） 取适当空间坐标系，求平面 \(BPC\) 与平面 \(APC\) 的法向量，代入二面角公式可得二面角的余弦值为 \(\dfrac{\sqrt{21}}{7}\).
\end{solution}

\begin{problem}
在 \(\triangle ABC\) 中，角 \(A,B,C\) 所对的边分别为 \(a,b,c\)，且 \(c=5\).
\begin{enumerate}
\item 若 \(\dfrac{a}{4b}=\dfrac{\sin B}{\sin A}\)，\(C=\dfrac{\pi}{2}\)，求 \(a\)；
\item 若 \(ab=20\)，求 \(\triangle ABC\) 面积的最大值.
\end{enumerate}
\end{problem}

\begin{solution}
（1） 由正弦定理，\(\dfrac{\sin B}{\sin A}=\dfrac{b}{a}\)，所以 \(\dfrac{a}{4b}=\dfrac{b}{a}\)，得 \(a=2b\). 又 \(C=\dfrac{\pi}{2}\)，故 \(a^2+b^2=c^2=25\). 解得 \(b=\sqrt5\)，\(a=2\sqrt5\).

（2） 面积 \(S=\dfrac12ab\sin C=10\sin C\). 由余弦定理 \(c^2=a^2+b^2-2ab\cos C\) 与 \(ab=20,c=5\) 得 \(\cos C=\dfrac{a^2+b^2-25}{40}\). 在约束下使 \(\sin C\) 最大，等价于使 \(\cos C\) 合法且尽量接近 \(0\). 由条件可得最大面积为 \(\dfrac{5\sqrt{55}}4\).
\end{solution}

\begin{problem}
甲、乙是两个体育社团的小组，给出两组组员身高的茎叶图（单位：厘米），以身高的百位数和十位数作为“茎”排列在中间、个位数作为“叶”分列在两边.

\begin{center}
\begin{tabular}{c|c|c}
甲 & 茎 & 乙\\
\hline
 & 15 & 9\\
7\quad 7\quad 5\quad 5\quad 4 & 16 & 0\quad 3\quad 5\quad 5\quad 6\quad 7\quad 8\quad 8\\
8\quad 5\quad 4\quad 3\quad 2\quad 2 & 17 & 2\\
3 & 18 &
\end{tabular}
\end{center}
\begin{enumerate}
\item 分别求甲、乙两组组员身高的第 \(60\) 百分位数；
\item 从甲、乙两组各选取一个组员，求两人身高均在 \(170\) 厘米以上的概率；
\item 为使两组人数相同，从甲组中调派一个队员到乙组. 是否存在甲组的一个组员，将他调派至乙组后，甲、乙两组的平均身高都增大？
\end{enumerate}
\end{problem}

\begin{solution}
（1） 按茎叶图从小到大排列，甲组第 \(60\) 百分位数为 \(173\)，乙组第 \(60\) 百分位数为 \(166.5\).

（2） 由茎叶图统计，甲组选到 \(170\) 厘米以上的可能数与乙组选到 \(170\) 厘米以上的可能数相乘，再除以总选法数，得概率为 \(\dfrac7{120}\).

（3） 存在. 调派身高介于两组原平均数之间的甲组成员，可使甲组去掉该成员后平均身高增大，同时乙组加入该成员后平均身高也增大；由茎叶图可取身高 \(167\) 厘米的成员.
\end{solution}

\begin{problem}
在平面直角坐标系中，已知曲线 \(\Gamma:\dfrac{x^2}{4}+y^2=1(y\ge0)\)，点 \(P,Q\) 分别为 \(\Gamma\) 上不同的两点，\(T(t,0)\).

\begin{enumerate}
\item 求 \(\Gamma\) 所在椭圆的离心率；
\item 若 \(T(1,0)\)，\(Q\) 在 \(y\) 轴上，且 \(T\) 到直线 \(PQ\) 的距离为 \(\dfrac{\sqrt5}{5}\)，求 \(P\) 的坐标；
\item 是否存在 \(t\)，使得 \(\triangle TPQ\) 是以 \(T\) 为直角顶点的等腰直角三角形？若存在，求 \(t\) 的取值范围；若不存在，请说明理由.
\end{enumerate}

\bitmapfigure[max width=0.42\linewidth,max totalheight=3.2cm]{img/2025/shanghai_spring/q20_fig1.png}

\end{problem}

\begin{solution}
（1） 椭圆标准方程为 \(\dfrac{x^2}{4}+\dfrac{y^2}{1}=1\)，故 \(a=2,b=1,c=\sqrt3\)，离心率 \(e=\dfrac{\sqrt3}{2}\).

（2） 设 \(Q=(0,1)\)，\(P=(x,y)\) 且 \(\dfrac{x^2}{4}+y^2=1\). 写出直线 \(PQ\) 方程，用点到直线距离公式代入 \(T(1,0)\)，并结合椭圆方程解得 \(P=(2,0)\).

（3） 若 \(\triangle TPQ\) 以 \(T\) 为直角顶点且为等腰直角三角形，则 \(\overrightarrow{TP}\perp\overrightarrow{TQ}\)，且 \(|TP|=|TQ|\). 将 \(P,Q\) 代入半椭圆参数方程并消元，可得存在的 \(t\) 的范围为 \(\left[-\dfrac65,\dfrac65\right]\).
\end{solution}

\begin{problem}
已知函数 \(y=f(x)\) 的定义域是 \(D\). 对于 \(t\in D\)，定义集合 \(S_f(t)=\{x\mid f(x)\ge f(t)\}\).
\begin{enumerate}
\item \(f(x)=\log_2x\)，求 \(S_f(16)\)；
\item 对于集合 \(A\)，若对任意 \(x\in A\) 都有 \(-x\in A\)，则称 \(A\) 是对称集. 若 \(D\) 是对称集，证明：“函数 \(y=f(x)\) 是偶函数”的充要条件是“对任意 \(t\in D\)，\(S_f(t)\) 是对称集”；
\item 若 \(x\in\mathbb R\)，\(f(x)=e^x-\dfrac12mx^2\). 求 \(m\) 的取值范围，使得对于任意 \(t_1<t_2\in D\)，都有 \(S_f(t_2)\subseteq S_f(t_1)\).
\end{enumerate}
\end{problem}

\begin{solution}
（1） 因为 \(\log_2x\ge\log_216=4\)，所以 \(x\ge16\)，故 \(S_f(16)=[16,+\infty)\).

（2） 若 \(f\) 为偶函数，则 \(f(x)=f(-x)\). 当 \(x\in S_f(t)\) 时，\(f(x)\ge f(t)\)，于是 \(f(-x)=f(x)\ge f(t)\)，故 \(-x\in S_f(t)\)，即 \(S_f(t)\) 为对称集. 反之，任取 \(t\in D\). 因为 \(t\in S_f(t)\)，由对称性得 \(-t\in S_f(t)\)，所以 \(f(-t)\ge f(t)\). 再把 \(t\) 换成 \(-t\)，得 \(f(t)\ge f(-t)\). 因此 \(f(t)=f(-t)\)，函数为偶函数.

（3） 条件等价于 \(f(x)\) 在 \(\mathbb R\) 上单调不减. 因为 \(f'(x)=e^x-mx\)，需对任意 \(x\) 有 \(e^x-mx\ge0\). 当 \(m<0\) 时取 \(x\to-\infty\) 不成立；当 \(m\ge0\) 时需 \(m\le \inf_{x>0}\dfrac{e^x}{x}=e\). 故 \(m\in[0,e]\).
\end{solution}
